%If you are presenting work which has been previously published, acknowledge this here.
% ***************************************************
% How to introduce a previously published chapter
% ***************************************************
%This is an example of how you might introduce a chapter that has been published previously. 
%\cleartoevenpage
%\pagestyle{empty}	
%Use this command (above) to suppress the header from the preceding chapter.

% ***************************************************
% Example of an internal chapter
% ***************************************************
%This is an internal chapter of the thesis.
%If you have a long title, you can supply an abbreviated version to print in the Table of Contents using the optional argument to the \chapter command.
\chapter[Abbreviated title]{Deep Learning}
\label{Chap:label}	%CREATE YOUR OWN LABEL.
\pagestyle{headings}

% ********* Enter your text below this line: ********
To provide context for this project, a high-level overview of deep learning will be provided. 

% ***************************************************


% ********* Enter your text below this line: ********
\section{Deep Learning Model Architectures}

There are a large variety of deep learning models, each designed for different purposes and with
different strengths and weaknesses. Given how computationally expensive training and fine-tuning models is,
ensuring that an appropriate architecture is used (indirectly chosen when selecting a pre-trained model)
is an important decision key for this project's success. In this section, we will discuss and compare common
deep learning model architectures at a high-level in order to determine an appropriate selection for this project.

\section{Simple Neural Network}

Here, we will outline a basic neural network model to familiarise ourselves with the terminology that will be used
later in the report, as well as the project moving forward.

\newpage

\begin{figure}
    \caption{A Simple Neural Network}
    \label{samplenetwork}
    \begin{center}
        \begin{neuralnetwork}[height=5]
            \newcommand{\x}[2]{$x_#2$}
            \newcommand{\y}[2]{$\hat{y}_#2$}
            \newcommand{\hfirst}[2]{\small $h^{(1)}_#2$}
            \newcommand{\hsecond}[2]{\small $h^{(2)}_#2$}
            \inputlayer[count=2, bias=false, title=Input\\layer, text=\x]
            \hiddenlayer[count=4, bias=false, title=Hidden\\layer 1, text=\hfirst] \linklayers
            \outputlayer[count=1, title=Output\\layer, text=\y] \linklayers
        \end{neuralnetwork}
    \end{center}
\end{figure}

Figure \ref{samplenetwork} shows a simple neural network. A neural network consists of an input layer, which is a vector. 
Transforming the raw data (e.g. images, natural language, Java documentation, etc.,) into a vector is known
as \emph{encoding}. How encoding is performed differs for all models, and is an consequential decision when designing a model.
Common methods for encoding include: integer encoding, one-hot encoding, and word embeddings.
Word embeddings are of particular importance as they are used by CodeBERT, the pre-trained model
that will be used for this project.

\begin{figure}
    \caption{Diagram of a singular neuron}
    \begin{center}
            \includegraphics[scale=0.75]{Images/neuron_diagram.png}
    \end{center}
    \label{neurondiagram}
\end{figure}

\newpage

Each neuron sums all the values of the neurons connected to it (generally, this is \emph{all} the previous neurons.).
The edges between neurons have \emph{weights}, which are multiplied against the value before being added to the neuron. 
Furthermore, each neuron has a \emph{bias} which is added to the neuron after each calculation. Figure \ref{neurondiagram} shows this graphically.

When this sum is generated, it is then passed into an \emph{activation function}.
The activation function stabalises the value of each neuron, ensuring they do not drop to zero or grow exponentially.
These two issues are commonly referred to as
the \emph{vanishing gradient} and \emph{exploding gradient} problems, respectively. 
The choice of activation function is important. There is a trade-off between the computational cost of activiation
functions, and how effective they are. Common activation functions include: the sigmoid function,
hyberbolic tangent function, and gaussian function.

\subsection{Recurrent Neural Networks}

Processing language is difficult. The meaning of words is not just determined by itself, but also in what
context the word is being used in. That is, the meaning of a word is not just determined by itself,
but also by the words that surround it.
The simple neural network introduced above would not be able to handle such a situation as it has no ``memory''
of previous inputs, i.e. it lacks contextual information.

Recurrent Neural Networks are one methodology for neural networks to have ``memory'' 
and thereby consider what context an input is being used in.
Essentially, the state of an RNN at any input is determined by all previous inputs.
Having ``memory'' of previous inputs allows RNNs to perform complex tasks on sequence data such as: voice recognition
computer vision or natural language tasks such as translation. % YOU NEED A CITATION HERE :)

While RNNs have historically proved useful for NLP tasks \cite{DBLP:journals/corr/BahdanauCB14}, they have
numerous downsides. Perhaps the key issue with RNNs is that they are notoriously difficult to train,
due to the \emph{vanishing gradient} and \emph{exploding gradient} problems. Difficulty here meaning both 
that they are computationally expensive to train, and can often produce inaccurate results for long sequences of inputs.
Most solutions to these issues slow training considerably \cite{DBLP:journals/corr/abs-1211-5063}.
LSTM (Long short-term memory) networks are a class of RNN conceived
to address the vanishing gradient and exploding gradient problems present in simple RNNs
\cite{10.1162/neco.1997.9.8.1735}. 

\subsection{Long short-term Memory Networks}

Long short-term memory networks (LTSMs) are a subclass of RNN designed to address the
vanishing gradient and exploding gradient issues present in simple RNNs. Describing the precise methodology
in which this is done is outside the scope of this report. However, put simply, this is achieved with ``gating units''
which are a more sophisticated method for calculating the hidden-state of a cell
(Where the hidden-state of a cell is the ``memory'' the neural network uses to understand the context of an input \cite{salem2022recurrent}).

These gating units, in addition to using the activation function that simple RNNs use, have to calculate a variety
of additional values. These extra computations accumulate and make LTSMs far more computationally intensive to train than
their generic RNN counterparts. However, this is not without benefit. LTSMs are significantly more accurate for longer
sequences compared to generic RNNs \cite{sundermeyer2012lstm}.


\subsection{Transformers}

The goal of the transformer model, introduced by Google Brain in 2017, was to handle sequential data, just as
RNNs do, but without the training-time limitations . Ultimately, this goal
was achieved as transformer models were demonstrated to perform significantly better than RNNs while
simultaneously requiring less training time \cite{DBLP:journals/corr/VaswaniSPUJGKP17}.

Long training times for RNNs are, ultimately, caused by the fact that input must be processed sequentially.
Transformers, however, do not have this limitation.
They can process multiple inputs simultaneously and those inputs do not have to be in order.
It is because transformers are parallelisable that they have quicker training times \cite{DBLP:journals/corr/VaswaniSPUJGKP17}.

